\begin{abstract}
\noindent
Federated learning fits a shared model across sites that cannot pool
raw records. When the scientific target is \emph{which variables matter}
rather than prediction alone, the natural tool is a sparse regression such
as the Lasso, and every federated protocol must choose $E$, the number of
local passes each site takes between aggregation rounds. Larger $E$ buys
fewer communication rounds at the risk of client drift. We study this
trade-off for a synchronous federated coordinate-descent Lasso and report a
finding that changes how such algorithms should be read: the weighted
average of locally soft-thresholded coefficient vectors is, except in
degenerate cases, not sparse, and its active set is the union of the local
active sets rather than the active set of any Lasso fit. We prove this in
closed form for orthogonal designs, where we also show that the iteration
terminates after a single round and is therefore entirely insensitive to
$E$, and we characterise the fixed point of the general iteration as an
averaged-stationarity condition that does not coincide with stationarity of
the aggregated objective. Consequently, published comparisons of active-set
size or support agreement across $E$ measure an artefact of the aggregation
rule as much as client drift. We propose a one-round remedy, a server-side
soft-threshold whose level is chosen by federated validation using scalar
summaries only, and we compare both against consensus ADMM, which minimises
exactly the same aggregated objective. In a Monte Carlo study spanning
independent, correlated and heterogeneous site designs, plain averaging
selects a large fraction of all predictors at every $E$, the server-side
threshold recovers most of the lost selection accuracy at negligible
communication cost, and consensus ADMM dominates both when sites are alike, though not
when they differ, since exact optimisation of a mis-specified common
penalty is no advantage.
We also show that ranking schedules by communication rounds alone is
misleading, since total local computation moves in the opposite direction.
\end{abstract}

\noindent\textbf{Keywords:} federated learning; Lasso; coordinate descent;
variable selection; communication efficiency; ADMM; client drift.
